\documentclass[../../main.tex]{subfiles}% 注意这里的文件路径不能用 ./main.tex ,否则用latexmk编译子文件会报错
\graphicspath{{\subfix{./image/}}} % 指定图片目录,后续可以直接使用图片文件名
% 注意这里的文件路径不能用 ../../image/ ,否则用latexmk编译子文件会报错

% 例如:
% \begin{figure}[H]
% \centering
% \includegraphics[scale=0.3]{图.png}
% \caption{}
% \label{figure:图}
% \end{figure}
% 注意:上述\label{}一定要放在\caption{}之后,否则引用图片序号会只会显示??.

\begin{document}

\section{紧算子的定义和基本性质}

\begin{definition}
设$\mathscr{X},\mathscr{Y}$是Banach空间，$A:\mathscr{X}\to\mathscr{Y}$为线性算子。若$\overline{A(B_1)}$在$\mathscr{Y}$中是紧集，其中$B_1$是$\mathscr{X}$中的闭单位球，则称$A$为\textbf{紧算子}。
全体紧算子构成的集合记作$\mathfrak{C}(\mathscr{X},\mathscr{Y})$，当$\mathscr{X}=\mathscr{Y}$时记作$\mathfrak{C}(\mathscr{X})$。
\end{definition}
\begin{remark}
$A\in\mathfrak{C}(\mathscr{X},\mathscr{Y})$的充要条件是：对$\mathscr{X}$中任意有界集$B$，$\overline{A(B)}$在$\mathscr{Y}$中为紧集。
\end{remark}
\begin{remark}
$A\in\mathfrak{C}(\mathscr{X},\mathscr{Y})$的充要条件是：对$\mathscr{X}$中任意有界点列$\{x_n\}$，$\{Ax_n\}$存在收敛子列。
\end{remark}

\begin{proposition}\label{proposition:泛函分析--命题3.1.2}
紧算子具有如下基本性质：
\begin{enumerate}[(1)]
\item $\mathfrak{C}(\mathscr{X},\mathscr{Y})\subseteq\mathscr{L}(\mathscr{X},\mathscr{Y})$；
\item 若$A,B\in\mathfrak{C}(\mathscr{X},\mathscr{Y})$，$\alpha,\beta\in\mathbb{C}$，则$\alpha A+\beta B\in\mathfrak{C}(\mathscr{X},\mathscr{Y})$；
\item\label{proposition:泛函分析--命题3.1.2-3} $\mathfrak{C}(\mathscr{X},\mathscr{Y})$在$\mathscr{L}(\mathscr{X},\mathscr{Y})$中是闭子集；
\item 若$A\in\mathfrak{C}(\mathscr{X},\mathscr{Y})$，$\mathscr{X}_0\subseteq\mathscr{X}$为闭线性子空间，则$A|_{\mathscr{X}_0}\in\mathfrak{C}(\mathscr{X}_0,\mathscr{Y})$；
\item 若$A\in\mathfrak{C}(\mathscr{X},\mathscr{Y})$，则值域$R(A)$可分；
\item 若$A\in\mathscr{L}(\mathscr{X},\mathscr{Y})$，$B\in\mathscr{L}(\mathscr{Y},\mathscr{Z})$，且其中一个算子为紧算子，则$BA\in\mathfrak{C}(\mathscr{X},\mathscr{Z})$。
\end{enumerate}
\end{proposition}
\begin{proof}
\begin{enumerate}[(1)]
\item 映射$y\mapsto\|y\|$连续，故
\begin{align*}
M\triangleq\sup\limits_{x\in B_1}\|Ax\|=\max\limits_{y\in\overline{A(B_1)}}\|y\|<\infty,
\end{align*}
因此$\|Ax\|\leqslant M\|x\|,\ \forall x\in\mathscr{X}$，即$A\in\mathscr{L}(\mathscr{X},\mathscr{Y})$。

\item 紧集的线性组合的闭包仍为紧集，故$\alpha A+\beta B$为紧算子。

\item 设$T_n\in\mathfrak{C}(\mathscr{X},\mathscr{Y})$，$\|T_n-T\|\to0\ (n\to\infty)$。
对任意$\varepsilon>0$，取$n\in\mathbb{N}$使得$\|T_n-T\|<\frac{\varepsilon}{2}$。
因$\overline{T_n(B_1)}$紧，存在有限$\frac{\varepsilon}{2}$-网$\{y_1,y_2,\cdots,y_m\}$，则
\begin{align*}
\overline{T(B_1)}\subseteq\bigcup\limits_{i=1}^{m}B(y_i,\varepsilon),
\end{align*}
即$\overline{T(B_1)}$有有限$\varepsilon$-网，故$\overline{T(B_1)}$为紧集，$T\in\mathfrak{C}(\mathscr{X},\mathscr{Y})$。

\item $\mathscr{X}_0$的有界集也是$\mathscr{X}$的有界集，紧算子限制后仍满足紧算子定义。

\item $R(A)=\bigcup\limits_{n=1}^{\infty}nA(B_1)$，$A(B_1)$的闭包紧，故$A(B_1)$可分，进而$R(A)$可分。

\item 连续线性算子将有界集映为有界集，将紧集映为紧集，故复合算子满足紧算子定义。
\end{enumerate}
\end{proof}

\begin{definition}[全连续算子]
设$A\in\mathscr{L}(\mathscr{X},\mathscr{Y})$，若对任意$x_n\rightharpoonup x$（弱收敛），均有$Ax_n\to Ax$（强收敛），则称$A$为\textbf{全连续算子}。
\end{definition}

\begin{proposition}
若$A\in\mathfrak{C}(\mathscr{X},\mathscr{Y})$，则$A$全连续；若$\mathscr{X}$是自反Banach空间且$A$全连续，则$A\in\mathfrak{C}(\mathscr{X},\mathscr{Y})$。
\end{proposition}
\begin{proof}
{\heiti 必要性：}设$x_n\rightharpoonup x$，反证。若存在$\varepsilon_0>0$及子列$\{x_{n_i}\}$使得$\|Ax_{n_i}-Ax\|\geqslant\varepsilon_0$。
由共鸣定理，$\{x_n\}$有界；由$A$紧，$\{Ax_{n_i}\}$存在收敛子列，不妨记为$Ax_{n_i}\to z$。
对任意$y^*\in\mathscr{Y}^*$，有
\begin{align*}
\langle y^*,Ax_{n_i}-Ax\rangle=\langle A^*y^*,x_{n_i}-x\rangle\to0,
\end{align*}
故$z=Ax$，与假设矛盾。

{\heiti 充分性:}由Eberlein-Smulian定理，$\mathscr{X}$中有界点列必有弱收敛子列；结合全连续性，该子列的像强收敛，故$A$为紧算子。

\end{proof}

\begin{theorem}
$T\in\mathfrak{C}(\mathscr{X},\mathscr{Y})$当且仅当$T^*\in\mathfrak{C}(\mathscr{Y}^*,\mathscr{X}^*)$。
\end{theorem}
\begin{proof}
{\heiti 必要性:}设$\{y_n^*\}$为$\mathscr{Y}^*$的单位球中有界点列，定义
\begin{align*}
\varphi_n(y)\triangleq\langle y_n^*,y\rangle,\quad \forall y\in\overline{T(B_1)}.
\end{align*}
则$\{\varphi_n\}\subseteq C(\overline{T(B_1)})$，满足
\begin{align*}
|\varphi_n(y)|\leqslant\|y_n^*\|\cdot\|y\|\leqslant\|T\|,\quad
|\varphi_n(y)-\varphi_n(z)|\leqslant\|y-z\|,\quad \forall y,z\in\overline{T(B_1)}.
\end{align*}
由Arzelà-Ascoli定理，$\{\varphi_n\}$存在一致收敛子列，对应$\{T^*y_n^*\}$存在收敛子列，故$T^*\in\mathfrak{C}(\mathscr{Y}^*,\mathscr{X}^*)$。

{\heiti 充分性:}由必要性得$T^{**}\in\mathfrak{C}(\mathscr{X}^{**},\mathscr{Y}^{**})$，而$T=T^{**}|_{\mathscr{X}}$，结合紧算子的限制性质，得$T\in\mathfrak{C}(\mathscr{X},\mathscr{Y})$。

\end{proof}

\begin{example}
设$\Omega\subseteq\mathbb{R}^n$为有界闭集，$K\in C(\Omega\times\Omega)$，$\mathscr{X}=\mathscr{Y}=C(\Omega)$，定义积分算子
\begin{align*}
T:u\mapsto\int_{\Omega}K(x,y)u(y)\mathrm{d}y,\quad \forall u\in C(\Omega),
\end{align*}
则$T\in\mathfrak{C}(\mathscr{X})$。
\begin{proof}
记$M\triangleq\max\limits_{x,y\in\Omega}|K(x,y)|$，则$\|Tu\|\leqslant M\|u\|\mathrm{mes}(\Omega)$。
由$K(x,y)$一致连续，对任意$\varepsilon>0$，存在$\delta>0$，当$|x-x'|<\delta$时，$|K(x,y)-K(x',y)|<\varepsilon$，故
\begin{align*}
|(Tu)(x)-(Tu)(x')|\leqslant\int_{\Omega}|K(x,y)-K(x',y)|\cdot|u(y)|\mathrm{d}y\leqslant\varepsilon\|u\|\mathrm{mes}(\Omega).
\end{align*}
由Arzelà-Ascoli定理，$\overline{T(B_1)}$为紧集，即$T\in\mathfrak{C}(\mathscr{X})$。

\end{proof}
\end{example}

\begin{example}
设$\Omega\subseteq\mathbb{R}^n$为有界开区域，$A\in\mathscr{L}(H_0^1(\Omega))$满足$\|Au\|_{H_0^1(\Omega)}\leqslant C\|u\|_{L^2(\Omega)}$，则$A\in\mathfrak{C}(H_0^1(\Omega))$。
\begin{proof}
由Rellich定理，嵌入算子$\iota:H_0^1(\Omega)\to L^2(\Omega)$为紧算子；$A:L^2(\Omega)\to H_0^1(\Omega)$连续，由紧算子的复合性质得证。

\end{proof}
\end{example}

\begin{definition}
设$T\in\mathscr{L}(\mathscr{X},\mathscr{Y})$，若$\dim R(T)<\infty$，则称$T$为\textbf{有限秩算子}，全体有限秩算子的集合记作$F(\mathscr{X},\mathscr{Y})$，且$F(\mathscr{X},\mathscr{Y})\subseteq\mathfrak{C}(\mathscr{X},\mathscr{Y})$。
\end{definition}

\begin{definition}
设$f\in\mathscr{X}^*$，$y\in\mathscr{Y}$，定义算子$y\otimes f$：
\begin{align*}
(y\otimes f)(x)=\langle f,x\rangle y,\quad \forall x\in\mathscr{X},
\end{align*}
称该算子为\textbf{秩1算子}。
\end{definition}

\begin{theorem}
$T\in F(\mathscr{X},\mathscr{Y})$的充要条件是：存在$y_i\in\mathscr{Y}$，$f_i\in\mathscr{X}^*\ (i=1,2,\cdots,n)$，使得
\begin{align*}
T=\sum\limits_{i=1}^{n}y_i\otimes f_i.
\end{align*}
\end{theorem}

\begin{proof}
充分性：$R(T)=\mathrm{span}\{y_1,y_2,\cdots,y_n\}$，故$\dim R(T)<\infty$，$T\in F(\mathscr{X},\mathscr{Y})$。

必要性：取$R(T)$的一组基$\{y_1,y_2,\cdots,y_n\}$，则对任意$x\in\mathscr{X}$，存在唯一线性泛函$l_i(x)$使得$Tx=\sum\limits_{i=1}^{n}l_i(x)y_i$。
因$\dim R(T)<\infty$，$\sum\limits_{i=1}^{n}|l_i(x)|$与$\|Tx\|$为等价范数，故存在$M>0$使得
\begin{align*}
\sum\limits_{i=1}^{n}|l_i(x)|\leqslant M\|Tx\|\leqslant M\|T\|\cdot\|x\|,\quad \forall x\in\mathscr{X},
\end{align*}
即$l_i\in\mathscr{X}^*$。记$f_i=l_i$，则$Tx=\sum\limits_{i=1}^{n}\langle f_i,x\rangle y_i=\left(\sum\limits_{i=1}^{n}y_i\otimes f_i\right)(x)$。

\end{proof}

\begin{remark}
\begin{enumerate}[(1)]
\item 若$\mathscr{X}$为Hilbert空间，则$\overline{F(\mathscr{X})}=\mathfrak{C}(\mathscr{X})$。
任取$T\in\mathfrak{C}(\mathscr{X})$，$\overline{T(B_1)}$紧，对任意$\varepsilon>0$，存在有限$\frac{\varepsilon}{2}$-网$\{y_1,\cdots,y_n\}$。令$E_\varepsilon=\mathrm{span}\{y_1,\cdots,y_n\}$，$P_\varepsilon$为正交投影，则$P_\varepsilon T\in F(\mathscr{X})$，且$\|T-P_\varepsilon T\|\leqslant\varepsilon$。
\item 若$\mathscr{X}$为一般Banach空间，结合紧算子值域可分，仅需讨论可分Banach空间。
\end{enumerate}
\end{remark}

\begin{definition}
设$\mathscr{X}$为可分Banach空间，若点列$\{e_n\}_{n=1}^{\infty}\subseteq\mathscr{X}$满足：对任意$x\in\mathscr{X}$，存在唯一数列$\{C_n(x)\}$使得
\begin{align*}
x=\lim\limits_{N\to\infty}\sum\limits_{n=1}^{N}C_n(x)e_n,
\end{align*}
则称$\{e_n\}$为$\mathscr{X}$的一组$\mathbf{Schauder}$\textbf{基}，$C_n(x)$为$\mathscr{X}$上的线性函数。
\end{definition}

\begin{lemma}
Schauder基对应的系数函数$C_n(x)\ (\forall n\in\mathbb{N})$是$\mathscr{X}$上的连续线性泛函。
\end{lemma}

\begin{proof}
在$\mathscr{X}$上引入新范数：$\square x\square=\sup\limits_{N\in\mathbb{N}}\left\|\sum\limits_{n=1}^{N}C_n(x)e_n\right\|$。
易验证$\mathscr{X}$关于$\square\cdot\square$完备，且$\|x\|=\lim\limits_{N\to\infty}\left\|\sum\limits_{n=1}^{N}C_n(x)e_n\right\|\leqslant\square x\square$。
由等价范数定理，存在$M_1>0$使得$\square x\square\leqslant M_1\|x\|$。
记$S_Nx=\sum\limits_{n=1}^{N}C_n(x)e_n$，则
\begin{align*}
\|C_n(x)e_n\|=\|S_nx-S_{n-1}x\|\leqslant2\square x\square\leqslant2M_1\|x\|,\quad \forall x\in\mathscr{X},
\end{align*}
故$|C_n(x)|\leqslant2M_1\|e_n\|^{-1}\|x\|$，即$C_n(x)$连续。

\end{proof}

\begin{theorem}
若可分Banach空间$\mathscr{X}$具有Schauder基，则$\overline{F(\mathscr{X})}=\mathfrak{C}(\mathscr{X})$。
\end{theorem}
\begin{proof}
对任意$N\in\mathbb{N}$，记$S_Nx=\sum\limits_{n=1}^{N}C_n(x)e_n$，$R_N=I-S_N$。
由共鸣定理，存在$M>0$使得$\|S_N\|\leqslant M$，故$\|R_N\|\leqslant1+M$。

任取$T\in\mathfrak{C}(\mathscr{X})$，对任意$\varepsilon>0$，因$\overline{T(B_1)}$紧，存在有限$\frac{\varepsilon}{3(M+1)}$-网$\{y_1,\cdots,y_m\}$，即
\begin{align}
\|Tx-y_i\|<\frac{\varepsilon}{3(M+1)},\quad \forall x\in B_1. \label{eq:compact_approx1}
\end{align}
由Schauder基定义，存在$N\in\mathbb{N}$使得
\begin{align}
\|y_j-S_Ny_j\|<\frac{\varepsilon}{3},\quad j=1,2,\cdots,m. \label{eq:compact_approx2}
\end{align}
结合\eqref{eq:compact_approx1}与$\|S_N\|\leqslant M$，得
\begin{align}
\|S_N(Tx)-S_Ny_i\|<\frac{M}{3(M+1)}\varepsilon. \label{eq:compact_approx3}
\end{align}
联立\eqref{eq:compact_approx1},\eqref{eq:compact_approx2},\eqref{eq:compact_approx3}，有
\begin{align*}
\|Tx-(S_NT)x\|<\varepsilon,\quad \forall x\in B_1,
\end{align*}
取$T_\varepsilon=S_NT\in F(\mathscr{X})$，则$\|T-T_\varepsilon\|\leqslant\varepsilon$，故$\overline{F(\mathscr{X})}=\mathfrak{C}(\mathscr{X})$。

\end{proof}

\begin{example}
设$\mathscr{X}$是一个无穷维$B$空间，求证：若$A\in\mathfrak{C}(\mathscr{X})$，则$A$没有有界逆。
\end{example}

\begin{example}
设$\mathscr{X}$是一个$B$空间，$A\in\mathscr{L}(\mathscr{X})$满足
\begin{align*}
\|Ax\|\geqslant\alpha\|x\|,\quad \forall x\in\mathscr{X},
\end{align*}
其中$\alpha$是正常数。求证：$A\in\mathfrak{C}(\mathscr{X})$的充要条件是$\mathscr{X}$是有穷维的。
\end{example}

\begin{example}
设$\mathscr{X},\mathscr{Y}$是$B$空间，$A\in\mathscr{L}(\mathscr{X},\mathscr{Y})$，$K\in\mathfrak{C}(\mathscr{X},\mathscr{Y})$，如果$R(A)\subseteq R(K)$，求证：$A\in\mathfrak{C}(\mathscr{X},\mathscr{Y})$。
\end{example}

\begin{example}
设$H$是Hilbert空间，$A:H\to H$是紧算子，又设$x_n\rightharpoonup x_0$，$y_n\rightharpoonup y_0$，求证：
\begin{align*}
(x_n,Ay_n)\to(x_0,Ay_0)\quad (n\to\infty).
\end{align*}
\end{example}

\begin{example}
设$\mathscr{X},\mathscr{Y}$是$B$空间，$A\in\mathscr{L}(\mathscr{X},\mathscr{Y})$，如果$R(A)$闭且$\dim R(A)=\infty$，求证：$A\notin\mathfrak{C}(\mathscr{X},\mathscr{Y})$。
\end{example}

\begin{example}
设$\omega_n\in\mathbb{K}$，$\omega_n\to0\ (n\to\infty)$，求证：映射
\begin{align*}
T:\{\xi_n\}\mapsto\{\omega_n\xi_n\}\quad (\forall\{\xi_n\}\in l^p)
\end{align*}
是$l^p(p\geqslant1)$上的紧算子。
\end{example}

\begin{example}
设$\Omega\subseteq\mathbb{R}^n$是一个可测集，又设$f$是$\Omega$上的有界可测函数，求证：$F:x(t)\mapsto f(t)x(t)$是$L^2(\Omega)$上的紧算子，当且仅当$f=0\ (\mathrm{a.e.}\text{于}\ \Omega)$。
\end{example}

\begin{example}
设$\Omega\subseteq\mathbb{R}^n$是一个可测集，又设$K\in L^2(\Omega\times\Omega)$，求证：
\begin{align*}
A:u(x)\mapsto\int_{\Omega}K(x,y)u(y)\mathrm{d}y\quad (\forall u\in L^2(\Omega))
\end{align*}
是$L^2(\Omega)$上的紧算子。
\end{example}

\begin{example}
设$H$是Hilbert空间，$A\in\mathfrak{C}(H)$，$\{e_n\}$是$H$的正交规范集，求证：$\lim\limits_{n\to\infty}(Ae_n,e_n)=0$。
\end{example}

\begin{example}
设$\mathscr{X}$是$B$空间，$A\in\mathfrak{C}(\mathscr{X})$，$\mathscr{X}_0$是$\mathscr{X}$的闭子空间并使得$A(\mathscr{X}_0)\subseteq\mathscr{X}_0$，求证：映射
\begin{align*}
T:[x]\mapsto[Ax]
\end{align*}
是商空间$\mathscr{X}/\mathscr{X}_0$上的紧算子。
\end{example}

\begin{example}
设$\mathscr{X},\mathscr{Y},\mathscr{Z}$是$B$空间，$\mathscr{X}\subseteq\mathscr{Y}\subseteq\mathscr{Z}$，如果$\mathscr{X}\to\mathscr{Y}$的嵌入映射是紧的，$\mathscr{Y}\to\mathscr{Z}$的嵌入映射是连续的，求证：$\forall\varepsilon>0$，$\exists c(\varepsilon)>0$，使得
\begin{align*}
\|x\|_{\mathscr{Y}}\leqslant\varepsilon\|x\|_{\mathscr{X}}+c(\varepsilon)\|x\|_{\mathscr{Z}}\quad (\forall x\in\mathscr{X}).
\end{align*}
\end{example}









\end{document}